% ============================================================================
% SECCIÓN 2.6.1: SÍNTESIS DEL ANÁLISIS COMPARATIVO
% Generado por: Poseidón 🔱 (12 Nov 2025, 21:45 hrs)
% Contenido: Análisis de brechas (5) + Justificación aportación diferencial
% ============================================================================

% ANÁLISIS DEL CUADRO:
\subsection{Síntesis del Análisis Comparativo}
\label{subsec:sintesis_comparativa}

Como se observa en la \Cref{tab:comparativa_investigaciones}, las investigaciones previas han abordado la clasificación del comportamiento sedentario y la actividad física desde múltiples perspectivas metodológicas, incluyendo deep learning, machine learning clásico, lógica difusa, métodos estadísticos avanzados y estudios de validación de dispositivos wearables. Sin embargo, el análisis comparativo revela la existencia de cinco brechas metodológicas y conceptuales significativas que limitan la aplicabilidad clínica y la generalización de los sistemas de clasificación actuales:

% ----------------------------------------------------------------------------
% BRECHA 1: Ausencia de sistemas híbridos clustering → fuzzy + LOUO
% ----------------------------------------------------------------------------

\paragraph{Brecha 1: Ausencia de sistemas híbridos clustering no supervisado + lógica difusa con validación LOUO.}

A pesar de la existencia de un precedente metodológico \cite{Goncalves2021} que combina clustering K-Means con un sistema de inferencia difusa Mamdani, \textbf{no se identifica ningún estudio que haya validado este enfoque híbrido mediante protocolos de Leave-One-User-Out (LOUO)}. La evidencia presentada por \cite{Rehman2024LOSO}, \cite{Mathew2024LOSO} y \cite{Bassani2025DLHAR} demuestra consistentemente que la validación LOUO produce caídas significativas en el rendimiento en comparación con métodos de validación menos rigurosos (e.g., k-fold, división 70/30), evidenciando problemas de generalización inter-usuario. Esta brecha es particularmente crítica para sistemas destinados a aplicaciones clínicas donde la generalización a nuevos pacientes es un requisito fundamental.

% ----------------------------------------------------------------------------
% BRECHA 2: Subutilización de HRV para clasificación fisiológica de sedentarismo
% ----------------------------------------------------------------------------

\paragraph{Brecha 2: Subutilización de HRV para la clasificación fisiológica del comportamiento sedentario.}

La mayoría de los estudios de clasificación de sedentarismo se basan exclusivamente en señales de movimiento (acelerómetro, giroscopio), abordando el problema desde una perspectiva postural (sentado, acostado, de pie) en lugar de fisiológica. El hallazgo de \cite{Godkin2025Context}, que demuestra que la frecuencia cardíaca en reposo (RHR) durante comportamiento sedentario es significativamente diferente de la RHR durante el sueño, evidencia la existencia de contextos autonómicos distintos. Adicionalmente, \cite{Marino2024ARIC} muestran que mayor HRV se asocia con mejor salud cognitiva, independientemente del nivel de actividad física. Sin embargo, \textbf{no se identifica ningún estudio que haya utilizado HRV como variable de entrada para clasificar comportamiento sedentario}, a pesar de su potencial para diferenciar estados fisiológicos (sedentarismo activo vs. sedentarismo pasivo) que no son discernibles mediante acelerometría sola.

% ----------------------------------------------------------------------------
% BRECHA 3: Dilema interpretabilidad vs. precisión
% ----------------------------------------------------------------------------

\paragraph{Brecha 3: Persistencia del dilema interpretabilidad vs. precisión en sistemas de clasificación.}

Los enfoques de deep learning \cite{Khan2024Wearable, Bassani2025DLHAR} alcanzan precisiones elevadas (90-95\%) pero operan como cajas negras, dificultando su adopción en contextos clínicos donde la confianza del profesional de salud depende de la transparencia del proceso de decisión \cite{Bienefeld2023XAI}. Por el contrario, los sistemas de lógica difusa \cite{Capitoli2025FuzzyXAI, MarashiHosseini2023Dietary} ofrecen interpretabilidad total mediante reglas lingüísticas, pero su precisión reportada proviene de aplicaciones en diagnóstico clínico (con variables clínicas discretas) y no en datos de sensores wearables continuos. \textbf{No existe evidencia de sistemas de lógica difusa aplicados a datos de Apple Watch} (acelerómetro + fotopletismografía) \textbf{para la clasificación de sedentarismo}, dejando sin resolver si la interpretabilidad puede preservarse sin sacrificar precisión en este dominio específico.

% ----------------------------------------------------------------------------
% BRECHA 4: Validación en cohortes de vida libre multi-anual
% ----------------------------------------------------------------------------

\paragraph{Brecha 4: Escasez de estudios con datos de vida libre multi-anual en contexto BYOD.}

La mayoría de los estudios identificados utilizan protocolos de recolección de datos de corta duración (días o semanas) en entornos semi-controlados. \cite{Fuller2021Predicting} validaron el paradigma Bring-Your-Own-Device (BYOD) con Apple Watch y Fitbit, pero con un período de seguimiento no especificado. \textbf{No se identifica ningún estudio que haya utilizado datos de Apple Watch recolectados durante múltiples años} (seguimiento longitudinal multi-anual) \textbf{en condiciones de vida libre sin restricciones}, a pesar de que este paradigma refleja con mayor fidelidad el uso real de wearables de consumo por parte de la población general. Esta brecha limita la validez ecológica de los modelos desarrollados en entornos controlados.

% ----------------------------------------------------------------------------
% BRECHA 5: Validación en cohortes pequeñas con N≤10
% ----------------------------------------------------------------------------

\paragraph{Brecha 5: Ausencia de metodologías validadas para cohortes pequeñas ($N \leq 10$) con variabilidad inter-individual alta.}

Los estudios identificados con validación LOUO utilizan cohortes de tamaño moderado a grande: \cite{Fuller2021Predicting} ($N=33$), \cite{Mathew2024LOSO} ($N=22$), \cite{OGrady2024AppleWatch} ($N=39$), \cite{Marino2024ARIC} ($N=961$). El caso de \cite{Mathew2024LOSO} con $N=22$ demostró una caída dramática en F1-Score (0.896 → 0.584) al aplicar LOUO, atribuida en parte al tamaño de muestra limitado. Sin embargo, \textbf{no se identifica ningún estudio que haya desarrollado y validado un sistema de clasificación con $N=10$ o menos}, un escenario común en estudios piloto, investigaciones clínicas con poblaciones raras, o fases iniciales de desarrollo de dispositivos médicos. \cite{Bolger2013} han demostrado que diseños longitudinales intensivos pueden compensar parcialmente tamaños de muestra pequeños, pero esta estrategia no ha sido aplicada en el contexto de clasificación de sedentarismo con wearables.

% ----------------------------------------------------------------------------
% JUSTIFICACIÓN DE LA APORTACIÓN DIFERENCIAL
% ----------------------------------------------------------------------------

\subsection{Aportación Diferencial de la Investigación Actual}
\label{subsec:aportacion_diferencial}

La presente investigación aborda simultáneamente las cinco brechas identificadas mediante un diseño metodológico novedoso que integra:

\begin{enumerate}
    \item \textbf{Sistema híbrido clustering no supervisado → lógica difusa supervisada con validación LOUO:} Se utiliza clustering K-Means ($k=2$) para establecer una ground truth operativa a partir de datos sin etiquetar de vida libre, seguido de un sistema de inferencia difusa Mamdani de 4 variables (Steps, HR\textsubscript{rest}, RMSSD, DailyEE/BMR) con 48 reglas lingüísticas. Este sistema es el \textbf{primero en ser validado mediante LOUO (10-fold)}, con resultados de F1-Score global de 0.840 y F1-Score LOOU de 0.780$\pm$0.167, demostrando generalización inter-usuario.
    
    \item \textbf{Inclusión de HRV (RMSSD) como variable fisiológica para clasificación de sedentarismo:} Se incorpora explícitamente la variabilidad de frecuencia cardíaca (RMSSD) como una de las cuatro variables de entrada del sistema difuso, permitiendo capturar el estado autonómico durante comportamiento sedentario. El análisis de ablación demuestra que la eliminación de HRV produce una caída del 50\% en F1-Score (4 variables → 2 variables), evidenciando su rol crítico en el modelo no lineal a pesar de no ser significativa en análisis univariado (Mann-Whitney $p=0.562$). Este es el \textbf{primer sistema que utiliza HRV para clasificar sedentarismo desde una perspectiva fisiológica}.
    
    \item \textbf{Interpretabilidad 100\% mediante reglas lingüísticas:} El sistema de lógica difusa Mamdani permite la inspección completa de las 48 reglas de inferencia y las funciones de membresía de cada variable, proporcionando transparencia total en el proceso de decisión. Esta interpretabilidad es comparable a la reportada por \cite{Capitoli2025FuzzyXAI} para diagnóstico clínico, pero aplicada por primera vez a datos de sensores Apple Watch continuos.
    
    \item \textbf{Validación con datos de vida libre multi-anual en contexto BYOD:} Se utilizan datos de Apple Watch recolectados durante un período de seguimiento medio de 133.7 semanas (rango: 7-298 semanas, equivalente a 0.1-5.7 años) en condiciones de vida libre sin restricciones. Este diseño longitudinal intensivo \cite{Bolger2013} compensa el tamaño de muestra pequeño mediante la recolección de 9,185 días totales de datos (1,337 semanas válidas tras control de calidad).
    
    \item \textbf{Metodología validada para cohorte pequeña ($N=10$, 5F/5M):} Se demuestra que es posible desarrollar un sistema de clasificación con generalización inter-usuario aceptable (F1-LOUO=0.780) incluso con $N=10$, mediante la combinación de: (a) seguimiento longitudinal extenso (compensando $N$ pequeño con alta densidad temporal), (b) validación LOUO rigurosa (evitando data leakage), y (c) modelo parsimonioso de 4 variables (evitando sobreajuste). Esta metodología es transferible a estudios piloto y poblaciones clínicas raras.
\end{enumerate}

En síntesis, ninguna de las investigaciones previas identificadas ha abordado simultáneamente la generalización inter-usuario (LOUO), la interpretabilidad total (lógica difusa), la clasificación fisiológica (HRV), la validez ecológica (vida libre multi-anual) y la aplicabilidad a cohortes pequeñas ($N \leq 10$). La presente investigación integra estos cinco elementos en un sistema funcional validado, posicionándola como una contribución metodológica novedosa al campo de la monitorización de comportamiento sedentario con wearables de consumo.

% ============================================================================
% FIN SECCIÓN 2.6.1
% ============================================================================

